\documentclass[11pt]{amsart}

\usepackage[margin=1.15in]{geometry}
\usepackage{amssymb}
\usepackage{booktabs}
\usepackage{needspace}
\usepackage[hidelinks]{hyperref}

\newcommand{\Q}{\mathbb{Q}}
\newcommand{\Z}{\mathbb{Z}}
\newcommand{\R}{\mathbb{R}}
\newcommand{\cod}[1]{\mathcal{D}_{#1}^{-1}}
\DeclareMathOperator{\Tr}{Tr}
\DeclareMathOperator{\Gal}{Gal}

\theoremstyle{plain}
\newtheorem{theorem}{Theorem}[section]
\newtheorem{lemma}[theorem]{Lemma}
\newtheorem{proposition}[theorem]{Proposition}
\newtheorem{corollary}[theorem]{Corollary}
\theoremstyle{definition}
\newtheorem{definition}[theorem]{Definition}
\newtheorem{remark}[theorem]{Remark}
\newtheorem{example}[theorem]{Example}

\title{A parity lemma for the codifferent of real cyclotomic fields}
\date{Draft, \today}

\begin{document}

\begin{abstract}
Let $n = n_1 n_2 n_3$ with $n_1, n_2, n_3 > 1$ odd and pairwise coprime, and let
$F_m = \Q(\zeta_m)^+$. We prove a congruence modulo~$2$ between the traces of an element of
the codifferent of $F_n$ to the three subfields $F_{n_i n_j}$. It shows, without any use of
positivity, that no element of the codifferent has the products $a_i a_j$ of trace-one
elements $a_i \in F_{n_i}$ as its three traces. Consequently, if each of the fields
$F_{pq}, F_{pr}, F_{qr}$ (for distinct odd primes $p, q, r$) has only products as its
level-two elements and $\{3,5\} \not\subseteq \{p,q,r\}$, then the minimal polynomial of
$2\cos(2\pi/pqr)$ is not $2$-interlacing.
We record what is known about this rigidity hypothesis, how the pair $\{3,5\}$ escapes the
congruence, and an explicit family of non-products in $F_{3q}$.
\end{abstract}

\maketitle

\section{Setting}\label{sec:setting}

For $n \geq 3$ let $\Psi_n$ be the minimal polynomial of $\alpha_n = 2\cos(2\pi/n)$, of degree
$d = \varphi(n)/2$, and let $F_n = \Q(\alpha_n) = \Q(\zeta_n)^+$, so that $O_{F_n} = \Z[\alpha_n]$
and the codifferent is $\cod{F_n} = \Z[\alpha_n]^\vee = \Psi_n'(\alpha_n)^{-1}\Z[\alpha_n]$.
A polynomial $g \in \Z[x]$ of degree $d-1$ with leading coefficient $j \geq 1$ interlaces
$\Psi_n$ (its roots are real and lie strictly between consecutive roots of $\Psi_n$) if and only
if $\gamma = g(\alpha_n)/\Psi_n'(\alpha_n)$ is a totally positive element
of $\cod{F_n}$ with $\Tr_{F_n/\Q}\gamma = j$: this is \cite[Prop.~3.5]{Y} for monic $g$ and the
remark following its proof in general.

\begin{definition}\label{def:level}
A \emph{level-$j$ element} of a totally real field $F$ is a totally positive
$\gamma \in \cod{F}$ with $\Tr_{F/\Q}\gamma = j$. For $F = F_n$ we write $m(n)$ for the least
$j$ for which a level-$j$ element exists; thus $\Psi_n$ is $j$-interlacing if and only if
$F_n$ has a level-$j$ element.
\end{definition}

If $E \subseteq F$ then $\Tr_{F/E}$ maps level-$j$ elements of $F$ to level-$j$ elements of $E$:
total positivity is clear, and for $\gamma \in \cod{F}$ and $e \in O_E$ we have
$\Tr_{E/\Q}(e\Tr_{F/E}\gamma) = \Tr_{F/\Q}(e\gamma) \in \Z$. Hence $m$ can only grow from a
subfield to the field.

\begin{theorem}[{\cite[Thm.~1.2]{Y}}]\label{thm:yatsyna}
Let $n$ be squarefree, not a prime or twice a prime, with $n > 20$ and $n \neq 30$. Then
$\Psi_n$ is non-interlacing, i.e.\ $m(n) \geq 2$.
\end{theorem}

The proof in \cite{Y} rests on \cite[Cor.~2.2]{BF}: a definite lattice with an automorphism of
characteristic polynomial $\Phi_m$, $m$ squarefree, $m \neq p, 2p$ and $\varphi(m) > 8$, does
not represent~$2$. For odd squarefree $n$ with at least two prime factors,
Theorem~\ref{thm:yatsyna} applies to every $n \neq 15$; and $m(15) = 1$.

The natural next question is whether $m(pqr) \geq 3$ for distinct odd primes $p, q, r$ (with
$\{3,5\} \not\subseteq \{p,q,r\}$, see Section~\ref{sec:three-five}). Computations in this
repository show that $F_n$ has no level-$2$ element for $n = 231$, $273$, $357$, $385$, $399$,
$429$, the six such $n$ below $455$ (\path{results/higher_interlacing_cyclotomic.md}).
This note isolates a mod-$2$ obstruction which reduces part of the question to two primes.

\section{Arrays}\label{sec:arrays}

Let $n \geq 3$ be odd, $M_n = \Q(\zeta_n)$ with $\zeta_n = e^{2\pi i/n}$, and $F_n = M_n \cap \R$.
For $\delta \in M_n$ define the \emph{array}
\[
  c_\delta \colon \Z/n \to \Q, \qquad c_\delta(x) = \Tr_{M_n/\Q}(\delta\,\zeta_n^x).
\]

\begin{lemma}\label{lem:arrays}
\begin{enumerate}
\item The map $\delta \mapsto c_\delta$ is injective and $\Q$-linear.
\item If $\delta \in F_n$ then $c_\delta(-x) = c_\delta(x)$ and $c_\delta(0) = 2\Tr_{F_n/\Q}\delta$.
\item If $\delta \in \cod{F_n}$ then $c_\delta$ takes values in $\Z$.
\item $\sum_{x \in \Z/n} c_\delta(x) = 0$.
\end{enumerate}
\end{lemma}

\begin{proof}
(1) The trace form of $M_n/\Q$ is non-degenerate and the $\zeta_n^x$ span $M_n$.
(2) For real $\delta$, $c_\delta(-x)$ is the complex conjugate of $c_\delta(x)$, which is
rational; and $\Tr_{M_n/\Q}\delta = 2\Tr_{F_n/\Q}\delta$.
(3) The different $\mathcal{D}_{M_n} = \mathcal{D}_{F_n}\mathcal{D}_{M_n/F_n}$ is contained in
$\mathcal{D}_{F_n}O_{M_n}$, so $\cod{F_n} \subseteq \cod{M_n}$, and
$\Tr_{M_n/\Q}(\delta\zeta_n^x) \in \Z$ for $\delta \in \cod{M_n}$ since $\zeta_n^x \in O_{M_n}$.
(4) $\sum_x \zeta_n^x = 0$.
\end{proof}

\begin{remark}\label{rem:fourier}
Writing $\sigma_k(\zeta_n) = \zeta_n^k$, one has
$\sum_{x} c_\delta(x)\zeta_n^{-lx} = n\,\sigma_l(\delta)$ for $l \in (\Z/n)^\times$ and $0$ for
non-units $l$, so $\delta \in F_n$ is totally positive if and only if this discrete Fourier
transform is positive on the units. Also $c_\delta(x - y) = \Tr_{M_n/\Q}(\delta\,\zeta_n^x
\overline{\zeta_n^y})$ is the Gram matrix of the ideal lattice
$(\Z[\zeta_n], \Tr_{M_n/\Q}(\delta u \bar v))$ in the basis of powers of $\zeta_n$; cf.\ \cite[Lemma~1.5]{BF}. Neither fact is used below.
\end{remark}

\subsection*{Blocks} Now let $n = n_1 \cdots n_k$ with $n_1, \dots, n_k > 1$ odd and pairwise
coprime. Identify $\Z/n = \prod_i \Z/n_i$ by $x \mapsto (x \bmod n_i)_i$, and let $e_i \in \Z/n$
be the idempotent with $e_i \equiv 1 \pmod{n_i}$ and $e_i \equiv 0 \pmod{n_j}$ for $j \neq i$.
Then $\zeta_{(i)} = \zeta_n^{e_i}$ is a primitive $n_i$-th root of unity and
$\zeta_n^x = \prod_i \zeta_{(i)}^{x_i}$. For a non-empty $S \subseteq \{1, \dots, k\}$ put
$n_S = \prod_{i \in S} n_i$ and $\zeta_S = \prod_{i \in S}\zeta_{(i)}$, a primitive $n_S$-th root
of unity generating $M_{n_S} \subseteq M_n$. \emph{Arrays of elements of $M_{n_S}$ are taken with
respect to $\zeta_S$}, as functions on $\prod_{i \in S}\Z/n_i$. (Another choice of primitive root
composes the array with multiplication by a unit; nothing below depends on it.)

\begin{lemma}\label{lem:blocks}
Let $\delta \in M_n$.
\begin{enumerate}
\item \textup{(Line sums.)} For every $i$ and every $x$,
  $\sum_{t \in \Z/n_i} c_\delta(x + t e_i) = 0$; that is, $c_\delta$ sums to zero along every line
  parallel to a coordinate axis.
\item \textup{(Traces are restrictions.)} For non-empty $S$,
  $c_{\Tr_{M_n/M_{n_S}}\delta}(x_S) = c_\delta(x_S, 0)$, where $(x_S, 0)$ has coordinates $x_i$
  for $i \in S$ and $0$ otherwise. If $\delta \in F_n$ then
  $\Tr_{M_n/M_{n_S}}\delta = \Tr_{F_n/F_{n_S}}\delta$.
\item \textup{(Products are tensors.)} Let $\{1, \dots, k\} = S \sqcup S'$. If
  $a \in M_{n_S}$ and $b \in M_{n_{S'}}$, then $c_{ab}(x_S, x_{S'}) = c_a(x_S)\,c_b(x_{S'})$.
\end{enumerate}
\end{lemma}

\begin{proof}
(1) The sum is $\Tr_{M_n/\Q}\bigl(\delta\zeta_n^x\sum_t\zeta_{(i)}^t\bigr) = 0$.
(2) For $w \in M_{n_S}$, $\Tr_{M_{n_S}/\Q}(w\Tr_{M_n/M_{n_S}}\delta) = \Tr_{M_n/\Q}(w\delta)$; take
$w = \zeta_S^{x_S} = \zeta_n^{(x_S,0)}$. Restriction $\Gal(M_n/M_{n_S}) \to \Gal(F_n/F_{n_S})$ is
injective, since an element of the kernel fixes $F_n$, hence is the identity or complex
conjugation, and complex conjugation moves $\zeta_S$ ($n_S \geq 3$); both groups have order
$\varphi(n)/\varphi(n_S)$.
(3) $M_n = M_{n_S} \otimes_\Q M_{n_{S'}}$ since $\varphi(n) = \varphi(n_S)\varphi(n_{S'})$, and the
trace of a tensor product of elements is the product of the traces.
\end{proof}

\begin{remark}[Composita]\label{rem:compositum}
Let $n = n'n''$ with $n', n''$ coprime and $a$, $b$ level-$j'$, level-$j''$ elements of
$F_{n'}$, $F_{n''}$. The fields $F_{n'}$, $F_{n''}$ have coprime discriminants, so
$O_{F_{n'}F_{n''}} = O_{F_{n'}} \otimes O_{F_{n''}}$ and $ab \in \cod{F_{n'}F_{n''}} \subseteq
\cod{F_n}$. As $[F_n : F_{n'}F_{n''}] = 2$, the element $ab$ has level $2j'j''$ in $F_n$; hence
$m(n'n'') \leq 2\,m(n')\,m(n'')$. We call these elements \emph{products}. By
Lemma~\ref{lem:blocks}(2),(3), if $\Tr_{F_{n''}/\Q} b = 1$ then
$\Tr_{F_n/F_{n'}}(ab) = 2a$, because both have the array $c_a(x')\,c_b(0) = 2c_a(x')$.
\end{remark}

\section{The parity lemma}\label{sec:parity}

Throughout this section $n = n_1 n_2 n_3$ with $n_i > 1$ odd and pairwise coprime, and
$A = \Z/n_1 \times \Z/n_2 \times \Z/n_3$. Let $\mathcal{A}$ be the set of functions
$c\colon A \to \Z$ with $c(-x) = c(x)$ and zero sums along all lines parallel to the axes. By
Lemmas~\ref{lem:arrays} and~\ref{lem:blocks}, $c_\delta \in \mathcal{A}$ for every
$\delta \in \cod{F_n}$.

A \emph{half-system} of $\Z/m$ ($m$ odd) is a subset $T \subseteq \Z/m \setminus \{0\}$ with
$\Z/m \setminus \{0\} = T \sqcup (-T)$; if $T$ is one, so is $-T$. Fix half-systems $T_i$ of
$\Z/n_i$.

\begin{lemma}\label{lem:identity}
For every $c \in \mathcal{A}$,
\[
  \sum_{(y,z) \in T_2 \times T_3} c(0,y,z) + \sum_{(x,z) \in T_1 \times (-T_3)} c(x,0,z)
  + \sum_{(x,y) \in T_1 \times T_2} c(x,y,0) \equiv 0 \pmod 2 .
\]
\end{lemma}

\begin{proof}
Call the three sums $X$, $Y$, $Z$. For signs $\varepsilon \in \{\pm\}^3$ let
$\Lambda(\varepsilon) = \sum c(x,y,z)$ over $(x,y,z) \in \varepsilon_1T_1 \times \varepsilon_2T_2
\times \varepsilon_3T_3$. Since $c$ is even, $\Lambda(-\varepsilon) = \Lambda(\varepsilon)$. Since
the line sums vanish and $\Z/n_1 = \{0\} \sqcup T_1 \sqcup (-T_1)$, summing over lines in the
first direction gives
$\Lambda(+,\varepsilon_2,\varepsilon_3) + \Lambda(-,\varepsilon_2,\varepsilon_3)
= -\sum_{\varepsilon_2T_2 \times \varepsilon_3T_3} c(0,y,z)$, and similarly in the other two
directions. Taking $(\varepsilon_2,\varepsilon_3) = (+,+)$ in the first direction,
$(\varepsilon_1,\varepsilon_3) = (+,-)$ in the second and $(\varepsilon_1,\varepsilon_2) = (+,+)$ in
the third, and using $\Lambda(-,+,+) = \Lambda(+,-,-)$,
\begin{align*}
  \Lambda(+,+,+) + \Lambda(+,-,-) &= -X,\\
  \Lambda(+,+,-) + \Lambda(+,-,-) &= -Y,\\
  \Lambda(+,+,+) + \Lambda(+,+,-) &= -Z.
\end{align*}
Adding, $X + Y + Z = -2\bigl(\Lambda(+,+,+) + \Lambda(+,-,-) + \Lambda(+,+,-)\bigr)$.
\end{proof}

\begin{definition}\label{def:tau}
Let $m_1, m_2 > 1$ be odd and coprime, and let $C\colon \Z/m_1 \times \Z/m_2 \to \Z$ be even,
with zero line sums and with even-valued axes $C(x,0)$, $C(0,y)$. Its \emph{corner parity} is
\[
  \tau(C) = \sum_{(x,y) \in T \times T'} C(x,y) \bmod 2
\]
for half-systems $T$, $T'$ of $\Z/m_1$, $\Z/m_2$. For $\beta \in \cod{F_{m_1m_2}}$ whose traces to
$F_{m_1}$ and $F_{m_2}$ have even-valued arrays we put $\tau(\beta) = \tau(c_\beta)$; by
Lemma~\ref{lem:blocks}(2) the axes of $c_\beta$ are the arrays of these two traces.
\end{definition}

\begin{lemma}\label{lem:tau-welldefined}
$\tau(C)$ does not depend on the half-systems.
\end{lemma}

\begin{proof}
Two half-systems differ by replacing some elements $t$ by $-t$. Replacing $t \in T'$ by $-t$
changes the sum by $\sum_{x \in T}\bigl(C(x,-t) - C(x,t)\bigr) \equiv \sum_{x \in T}\bigl(C(-x,t)
+ C(x,t)\bigr) = \sum_{x \neq 0} C(x,t) = -C(0,t) \equiv 0$, using $C(x,-t) = C(-x,t)$ and the
zero line sum through $(0,t)$. The same holds for $T$.
\end{proof}

\begin{theorem}[Parity lemma]\label{thm:parity}
Let $\delta \in \cod{F_n}$ and suppose that the arrays of $\Tr_{F_n/F_{n_i}}\delta$ are
even-valued for $i = 1, 2, 3$. Then
\[
  \tau\bigl(\Tr_{F_n/F_{n_2n_3}}\delta\bigr) + \tau\bigl(\Tr_{F_n/F_{n_1n_3}}\delta\bigr)
  + \tau\bigl(\Tr_{F_n/F_{n_1n_2}}\delta\bigr) \equiv 0 \pmod 2 .
\]
\end{theorem}

\begin{proof}
Let $c = c_\delta \in \mathcal{A}$. By Lemma~\ref{lem:blocks}(2) and transitivity of the trace, the
restriction of $c$ to the coordinate plane through $0$ orthogonal to the $i$-th axis is the array
of $\Tr_{F_n/F_{n_jn_k}}\delta$, and its axes are the arrays of the traces to $F_{n_j}$ and
$F_{n_k}$, which are even-valued by hypothesis. Each restriction is even with zero line sums, so
its corner parity is defined. The three sums of Lemma~\ref{lem:identity} compute these corner
parities with the half-systems $(T_2,T_3)$, $(T_1,-T_3)$ and $(T_1,T_2)$, and by
Lemma~\ref{lem:tau-welldefined} the choice does not matter.
\end{proof}

Positivity plays no role: Theorem~\ref{thm:parity} is a statement about $\cod{F_n}$ modulo~$2$.

\begin{lemma}\label{lem:product-tau}
Let $m_1, m_2 > 1$ be odd and coprime and let $a \in \cod{F_{m_1}}$, $b \in \cod{F_{m_2}}$ have
trace~$1$. Then the traces of $ab \in F_{m_1m_2}$ to $F_{m_1}$, $F_{m_2}$ are $2a$, $2b$, and
$\tau(ab) = 1$.
\end{lemma}

\begin{proof}
By Lemma~\ref{lem:blocks}(3), $c_{ab} = c_a \otimes c_b$, with axes $2c_a$, $2c_b$
(Remark~\ref{rem:compositum}). By Lemma~\ref{lem:arrays}, $c_a$ is integral and even,
$c_a(0) = 2$ and $\sum_x c_a(x) = 0$; so $\sum_{x \in T} c_a(x) = \frac12\sum_{x \neq 0}c_a(x) = -1$.
Likewise for $b$, and $\tau(ab) = \bigl(\sum_T c_a\bigr)\bigl(\sum_{T'} c_b\bigr) = 1$.
\end{proof}

\begin{corollary}[Product obstruction]\label{cor:products}
Let $a_i \in \cod{F_{n_i}}$ have trace~$1$ $(i = 1,2,3)$. There is no $\delta \in \cod{F_n}$ with
\[
  \Tr_{F_n/F_{n_2n_3}}\delta = a_2a_3,\qquad \Tr_{F_n/F_{n_1n_3}}\delta = a_1a_3,\qquad
  \Tr_{F_n/F_{n_1n_2}}\delta = a_1a_2 .
\]
\end{corollary}

\begin{proof}
For such a $\delta$, $\Tr_{F_n/F_{n_1}}\delta = \Tr_{F_{n_1n_2}/F_{n_1}}(a_1a_2) = 2a_1$, whose array
$2c_{a_1}$ is even; likewise for $n_2$, $n_3$. By Lemma~\ref{lem:product-tau} the three corner
parities are $1$, and $1 + 1 + 1$ is odd, contradicting Theorem~\ref{thm:parity}.
\end{proof}

\begin{remark}[The number of blocks]\label{rem:blocks}
The proof of Lemma~\ref{lem:identity} works for any number $k$ of blocks: walking from
$\varepsilon$ to $-\varepsilon$ by changing one sign at a time gives
$\Lambda(\varepsilon)\bigl(1 - (-1)^k\bigr) = $ an alternating sum of $k$ face sums. For odd $k$
this is a congruence modulo~$2$; for even $k$ it is an exact linear relation. For $k = 2$ the
relation reads $\sum_{T_1}c(x,0) = \sum_{T_2}c(0,y)$, and both sides equal $-c(0,0)/2$, so it
carries no information. This matches the fact that products of level-one elements are level-two
elements of $F_{pq}$ (Remark~\ref{rem:compositum}), while three-fold products are ruled out by
Corollary~\ref{cor:products}.
\end{remark}

\section{Three primes}\label{sec:three-primes}

\begin{definition}\label{def:rigid}
For coprime odd $m_1, m_2 > 1$, the field $F_{m_1m_2}$ is \emph{rigid at level two} if every
level-two element of $F_{m_1m_2}$ is a product $ab$ of level-one elements $a \in F_{m_1}$,
$b \in F_{m_2}$.
\end{definition}

\begin{theorem}\label{thm:three}
Let $n = n_1n_2n_3$ with $n_i > 1$ odd and pairwise coprime. Suppose that for each pair
$\{i,j\}$ the field $F_{n_in_j}$ has no level-one element and is rigid at level two. Then $F_n$
has no element of level $\leq 2$, i.e.\ $m(n) \geq 3$.
\end{theorem}

\begin{proof}
Let $\delta$ have level $j \leq 2$. Its traces to the fields $F_{n_in_j}$ have level $j$
(Section~\ref{sec:setting}), so $j = 2$ and they are products:
$\Tr_{F_n/F_{n_1n_2}}\delta = a_1a_2$, $\Tr_{F_n/F_{n_1n_3}}\delta = a_1'a_3$ and
$\Tr_{F_n/F_{n_2n_3}}\delta = a_2'a_3'$. Tracing further to $F_{n_1}$ and using
Remark~\ref{rem:compositum}, $2a_1 = \Tr_{F_n/F_{n_1}}\delta = 2a_1'$, so $a_1 = a_1'$; likewise
$a_2 = a_2'$ and $a_3 = a_3'$. This contradicts Corollary~\ref{cor:products}.
\end{proof}

\begin{corollary}\label{cor:primes}
Let $p, q, r$ be distinct odd primes with $\{3,5\} \not\subseteq \{p,q,r\}$. If $F_{pq}$, $F_{pr}$
and $F_{qr}$ are rigid at level two, then $m(pqr) \geq 3$: the minimal polynomial of
$2\cos(2\pi/pqr)$ is not $2$-interlacing.
\end{corollary}

\begin{proof}
Each product of two of the primes is squarefree, neither prime nor twice a prime, exceeds $20$
and differs from $30$, so it has no level-one element by Theorem~\ref{thm:yatsyna}. Apply
Theorem~\ref{thm:three}.
\end{proof}

Theorem~\ref{thm:three} also applies to composite blocks; for example, the blocks
$(15, p, q)$ need rigidity of $F_{15p}$, $F_{15q}$ and $F_{pq}$, with products built from the
level-one elements of $F_{15}$.

\section{Status and examples}\label{sec:status}

\subsection{Rigidity}\label{sec:rigidity}
Table~\ref{tab:rigid} is taken from \path{results/higher_interlacing_cyclotomic.md}, where
the level-two elements were enumerated by towers of subfields. Table~\ref{tab:census} sorts the
level-two elements of the smallest fields $F_{pq}$ by the type of their traces; it is reproduced
by \path{scripts/parity_table.py} (about $15$ seconds) and stored in
\path{results/parity_level2.json}.

\begin{table}[ht]
\centering
\begin{tabular}{lrrl}
\toprule
field & level-two elements & products & rigid \\
\midrule
$F_{77}$  & 285 & 15 & no \\
$F_{91}$  & 18  & 18 & yes \\
$F_{119}$ & 24  & 24 & yes \\
$F_{133}$ & 27  & 27 & yes \\
$F_{143}$ & 30  & 30 & yes \\
\bottomrule
\end{tabular}
\caption{Level-two elements of $F_{pq}$ found in the computations of $m(pqr)$ for $pqr < 455$.}
\label{tab:rigid}
\end{table}

\begin{table}[ht]
\centering
\begin{tabular}{lrrrrrr}
\toprule
 & & & & \multicolumn{2}{c}{non-products} & \\
\cmidrule(lr){5-6}
field & degree & level two & products & $\tau = 0$ & $\tau = 1$ & other \\
\midrule
$F_{15}$ & 4  & 50   & 2 & 12 & 12 & 24 \\
$F_{21}$ & 6  & 141  & 3 & 18 & 18 & 102 \\
$F_{33}$ & 10 & 935  & 5 & 20 & 30 & 880 \\
$F_{35}$ & 12 & 1470 & 6 & 36 & 84 & 1344 \\
$F_{39}$ & 12 & 2214 & 6 & 48 & 24 & 2136 \\
\bottomrule
\end{tabular}
\caption{Level-two elements $\beta$ of $F_{pq}$. The middle columns count those with
$\Tr_{F_{pq}/F_p}\beta = 2a$ and $\Tr_{F_{pq}/F_q}\beta = 2b$ for level-one elements $a$, $b$
(products, and non-products by corner parity); ``other'' counts the rest.}
\label{tab:census}
\end{table}

So no small $F_{pq}$ is rigid, and the rigid pairs $\{7,13\}$, $\{7,17\}$, $\{7,19\}$,
$\{11,13\}$ of Table~\ref{tab:rigid} contain no triangle: completing one needs a further field
such as $F_{221}$ or $F_{247}$ (degrees $96$ and $108$), for the triples $\{7,13,17\}$ and
$\{7,13,19\}$. \emph{At present
Corollary~\ref{cor:primes} proves no case that was not already settled by computation.}

\subsection{Gluing}\label{sec:gluing}
Theorem~\ref{thm:parity} also constrains the case where only two of the three fields are rigid.

\begin{proposition}\label{prop:glue}
Let $n = n_1n_2n_3$ be as in Section~\ref{sec:parity}, and let $\delta$ be a level-two element of
$F_n$ whose traces to $F_{n_1n_3}$ and $F_{n_2n_3}$ are products of level-one elements. Then
$\beta = \Tr_{F_n/F_{n_1n_2}}\delta$ is a level-two element with traces $2a_1$, $2a_2$ to
$F_{n_1}$, $F_{n_2}$ for level-one elements $a_1$, $a_2$; it is not a product, and $\tau(\beta) = 0$.
\end{proposition}

\begin{proof}
Write the two traces as $a_1a_3$ and $a_2a_3'$; tracing to $F_{n_3}$ gives $a_3 = a_3'$, and
tracing to $F_{n_i}$ gives $\Tr_{F_{n_1n_2}/F_{n_i}}\beta = 2a_i$. If $\beta$ were a product it
would be $a_1a_2$, contradicting Corollary~\ref{cor:products}. Theorem~\ref{thm:parity} with
Lemma~\ref{lem:product-tau} gives $1 + 1 + \tau(\beta) \equiv 0$.
\end{proof}

Hence, for $n = pqr$ with $F_{pr}$, $F_{qr}$ rigid, only the elements in the column
$\tau = 0$ of Table~\ref{tab:census} can occur as the trace to $F_{pq}$ of a level-two element.
For the six $n \leq 429$ of Section~\ref{sec:setting} the computations show that $F_n$ has no
level-two element at all. Parity alone does not explain this, since every field of
Table~\ref{tab:census} has non-products with $\tau = 0$; positivity must be used as well.

\subsection{The pair \texorpdfstring{$\{3,5\}$}{\{3,5\}}}\label{sec:three-five}
The field $F_{15}$ has eight level-one elements~$e$, and each $2e$ has an even-valued array, so
$\tau(2e) = 0$. For an odd prime $r \notin \{3,5\}$ and a level-one element $a_r$ of $F_r$, the
element $\delta = e\,a_r \in F_{15}F_r$ has level $2$ in $F_{15r}$ (Remark~\ref{rem:compositum}).
By Lemma~\ref{lem:blocks}, its traces to $F_{15}$, $F_{3r}$ and $F_{5r}$ are $2e$,
$(\Tr_{F_{15}/\Q}e)\,a_r = a_r$ and $(\Tr_{F_{15}/F_5}e)\,a_r$, with corner parities $0$, $1$,
$1$, as Proposition~\ref{prop:glue} requires. This is how $\{3,5\}$ passes the parity test, and
together with Theorem~\ref{thm:yatsyna} it gives $m(15r) = 2$ for every such $r$.

\needspace{8\baselineskip}
\subsection{Non-products in \texorpdfstring{$F_{3q}$}{F\_3q}}\label{sec:family}
\begin{proposition}\label{prop:family}
Let $q > 3$ be prime, $K = F_{3q}$, $E = F_q \subset K$, $a = (2 - \alpha_q)/q$ and
$\theta = (\zeta_3 - \zeta_3^{-1})(\zeta_q - \zeta_q^{-1}) \in K$. Then
$\gamma_\pm = a(1 \pm \theta^{-1})$ lie in $\cod{K}$, have trace $2$, do not lie in $E$, and are
totally positive if and only if $12\sin^2(\pi/q) > 1$, i.e.\ $q \in \{5, 7\}$.
\end{proposition}

\begin{proof}
Put $\alpha = \zeta_3\zeta_q + \zeta_3^{-1}\zeta_q^{-1}$, a conjugate of $\alpha_{3q}$, so that
$O_K = \Z[\alpha]$. Its conjugate over $E$ is
$\alpha' = \zeta_3^{-1}\zeta_q + \zeta_3\zeta_q^{-1}$, and $\alpha - \alpha' = \theta$. Since
$O_K = \Z[\alpha] \subseteq O_E[\alpha] \subseteq O_K$, the relative different is
$\mathcal{D}_{K/E} = (\alpha - \alpha')O_K = \theta O_K$. Next, $a$ is a level-one element of $E$: the prime $\mathfrak{p} = (2 - \alpha_q)$ of $E$ is
the only one above $q$, with $\mathfrak{p}^{(q-1)/2} = (q)$, and $E/\Q$ is tamely and totally
ramified at $q$ and unramified elsewhere, so $\mathcal{D}_E = \mathfrak{p}^{(q-3)/2}$ and $a$ generates
$\mathfrak{p}^{1-(q-1)/2} = \cod{E}$; it is totally positive, and
$\Tr_{E/\Q}a = (q - 1 - \Tr_{E/\Q}\alpha_q)/q = 1$ as $\Tr_{E/\Q}\alpha_q = -1$. Hence
$\cod{K} = \theta^{-1}\cod{E}O_K$ contains $a$ and $a/\theta$. The non-trivial automorphism of $K/E$
sends $\theta$ to $-\theta$, so $\Tr_{K/E}\theta^{-1} = 0$ and $\Tr_{K/\Q}\gamma_\pm =
\Tr_{E/\Q}(2a) = 2$; and $\gamma_\pm \notin E$. At a real place of $E$ with
$\alpha_q \mapsto 2\cos(2\pi k/q)$, $a$ is positive and $\theta^2 = (\alpha_3^2 - 4)(\alpha_q^2 - 4)
= 12\sin^2(2\pi k/q)$, and $\theta$ takes the two values $\pm\sqrt{\theta^2}$ at the two places above it; so
$\gamma_+$ and $\gamma_-$ are positive at both if and only if $\theta^2 > 1$ there. The least value of $|\sin(2\pi k/q)|$ over $k \in (\Z/q)^\times$ is $\sin(\pi/q)$,
and $12\sin^2(\pi/q)$ is $4.15$, $2.26$, $0.95$ for $q = 5, 7, 11$ and decreases in $q$.
\end{proof}

The elements $\gamma_\pm$ have traces $2a$ to $E$ and $2$ to $F_3 = \Q$; their corner parity
is $0$ for $q = 5$ and $1$ for $q = 7$. They account for some of the non-products of
Table~\ref{tab:census} in $F_{15}$ and $F_{21}$ only; the non-products in $F_{33}$, $F_{35}$,
$F_{39}$ are of other kinds.

\subsection{Open problems}\label{sec:open}
\begin{enumerate}
\item Prove that $F_{pq}$ is rigid at level two for all primes $p, q \geq p_0$. With
  Corollary~\ref{cor:primes} this gives $m(pqr) \geq 3$ for $p, q, r \geq p_0$.
\item For each of the finitely many non-rigid pairs expected, exclude the gluing of
  Proposition~\ref{prop:glue} uniformly in $r$, or find an exception like $\{3,5\}$.
\item For four or more primes one needs the minimal elements at the previous level, which are not
  products: products have level $4$ in $F_{pqr}$, while $m(231) = 3$.
\end{enumerate}

\subsection{Verification}\label{sec:verification}
\path{tests/test_parity.py} checks Lemma~\ref{lem:identity} on random elements of
$\cod{F_{105}}$ for three orderings of the blocks $(3,5,7)$, Lemma~\ref{lem:product-tau} in
$F_{35}$, the rows $F_{15}$ and $F_{21}$ of Table~\ref{tab:census}, and the positivity statement of
Proposition~\ref{prop:family} for $q = 5, 7, 11, 13$. Lemma~\ref{lem:identity} was also checked
on random elements of $\cod{F_{231}}$, and Proposition~\ref{prop:family} for $q \leq 29$. The
arrays are computed exactly as $c_\delta(x) = \Tr_{F_n/\Q}(\delta\,C_x(\alpha_n))$ with
$C_x(2\cos t) = 2\cos(xt)$ (\path{noninterlacing/parity.py}).

\begin{thebibliography}{BF}

\bibitem[BF]{BF}
E.~Bayer-Fluckiger, \emph{Definite unimodular lattices having an automorphism of given
characteristic polynomial}, Comment. Math. Helv. \textbf{59} (1984), 509--538.

\bibitem[Y]{Y}
P.~Yatsyna, \emph{A lower bound for the rank of a universal quadratic form with integer
coefficients in a totally real number field}, Comment. Math. Helv. \textbf{94} (2019),
221--239.

\end{thebibliography}

\end{document}
